\chapter{Metodologia Aplicada}

Este capítulo detalha as estratégias, as decisões de engenharia e os procedimentos experimentais adotados para a construção do protótipo EB-15. O foco central reside em traduzir os embasamentos teóricos sobre controle motor e Internet das Coisas (\textit{Internet of Things} - IoT) em implementações tangíveis, garantindo a viabilidade de uma plataforma de experimentação de baixo custo. A documentação dos métodos abrange desde a fabricação estrutural até a definição quantitativa dos testes aplicados ao conjunto final. Sendo assim, o registro sistemático de cada etapa possibilita avaliar as limitações tecnológicas da solução proposta.

Nesse contexto, a escolha minuciosa dos componentes e do fluxo de análise garante que o protótipo possa ser testado sob rigorosos critérios de estabilidade e precisão. Os ensaios delineados no planejamento fornecem métricas validadas de desempenho real para a posterior análise gráfica em comparação ao manipulador industrial UR-10. Deste modo, as intervenções metodológicas justificam as escolhas de \textit{hardware} em favor de soluções baseadas na otimização da lógica de \textit{software}. Portanto, a execução prática busca responder diretamente aos objetivos estabelecidos para a pesquisa acadêmica.

\section{Tipo e Caracterização da Pesquisa}

A pesquisa caracteriza-se pela sua natureza aplicada, mesclando metodologias de \textit{design} em engenharia de controle com observação experimental analítica. A construção do braço robótico serve como veículo primário para materializar o problema exposto na seção introdutória, validando a emulação de dinâmicas industriais usando orçamentos laboratoriais universitários \cite{siciliano2009}. Trata-se também de uma investigação avaliativa, visto que a performance do artefato acessível é prospectada à luz do referencial canônico de precisão ditado por equipamentos de automação modernos. Desse modo, o cruzamento de dados quantitativos embasa a viabilidade técnica da solução proposta.

Com base nisso, as intervenções da pesquisa ocorrem localmente e operam sob condições instrumentais limitadas, utilizando sensores óticos comerciais e paquímetros em laboratório. A investigação avança sistematicamente da integração isolada dos motores para a experimentação global, avaliando os seis eixos operando simultaneamente em rede remota. Por conseguinte, a obtenção de conclusões satisfatórias advém de um mapeamento estrito entre a intenção matemática de movimento prescrita no \textit{firmware} e a resposta temporal mensurada no mundo físico. Assim, garante-se uma base de dados sólida para as avaliações posteriores.

\section{Plataforma e Materiais Utilizados}

A seleção da arquitetura formadora do manipulador guiou-se pelo pilar do projeto \textit{open-source} do modelo EB-15, provendo um alicerce sólido. Essa base minimizou adaptações geométricas dispendiosas para alcançar os seis graus de liberdade necessários à simulação industrial. Em substituição ao metal usinado, empregou-se plástico modificado submetido a técnicas avançadas de impressão 3D de alta precisão. Consequentemente, essa estratégia possibilita validar o uso acadêmico de componentes sintéticos como alternativas funcionais para análises de cinemática e dinâmica.

Para tanto, os módulos robóticos basearam-se fortemente no arranjo mecânico engrenado com caixas de redução planetárias. Esse formato construtivo favorece a ampliação do torque, compensando as limitações energéticas típicas do ambiente experimental e dos motores de passo empregados. Todavia, como discutido na fundamentação formal por \citeonline{craig2018}, a elasticidade inerente ao plástico induz variações dimensionais não esperadas sob aquecimento e estresse contínuo. Dessa forma, as especificidades técnicas exigem algoritmos capazes de minimizar oscilações e incertezas ao longo das trajetórias operacionais.

\subsection{Fabricação e Estrutura do Protótipo EB-15}

A materialização estrutural e mecânica dos módulos deu-se via manufatura aditiva utilizando uma impressora 3D munida de extrusor de alta precisão. O material selecionado foi o filamento PETG, escolhido criteriosamente por seu notável equilíbrio entre resistência mecânica e flexibilidade. A manufatura empenhou cerca de um quilograma e meio de matéria-prima, processada em mais de cem horas operacionais para forjar as articulações e carcaças. Portanto, a integridade da estrutura física sustenta as avaliações motoras realizadas ao longo do desenvolvimento.

Por conseguinte, a mitigação de falhas elásticas induziu diretrizes paramétricas estritas para o fatiador tridimensional. Adotou-se o uso massivo de contornos reforçados delimitados por seis camadas sólidas e uma matriz de preenchimento tipo \textit{Gyroid} em áreas de pico de tensão. Adicionalmente, as engrenagens centrais receberam atenção prioritária para resistir ao torque dinâmico dos motores sob carga contínua. Sendo assim, a combinação metódica dessas parametrizações logrou absorver os choques inerciais oriundos da ausência de rampas de aceleração puramente mecânicas.

Todavia, a resistência inerente ao polímero impresso apresenta um risco de falha estrutural sob torques elevados. As caixas de redução planetárias, especificamente, podem ceder diante do esforço mecânico contínuo gerado pelos atuadores basilares. Diante dessa possibilidade de ruptura, a metodologia estipula a segmentação emergencial do protótipo em dois blocos de teste independentes. Desse modo, o módulo de movimentação (composto pelas três primeiras juntas) seria fisicamente isolado do módulo de direção (constituído pelas três últimas juntas). Portanto, essa estratégia de contingência permite a validação parcial dos componentes, garantindo a continuidade do cronograma experimental.

\subsection{Eixo de Atuadores e Eletrônica de Controle}

O esquema de acionamento dinâmico adotou atuadores de passo padrão NEMA 17 para traduzir o deslocamento elétrico em movimento físico. Para garantir o suprimento energético necessário à manutenção de corrente constante, aplicou-se circuitos tracionadores dedicados a cada junta basilar do arranjo de potência. Essa configuração mitiga o efeito indutivo indesejado que comprometeria a exatidão quantitativa durante operações sequenciais extensas. Por conseguinte, a estabilidade elétrica suporta as metodologias de rastreamento baseadas unicamente na predição de passos executados.

Entrementes, a inteligência motora segmenta-se em uma arquitetura distribuída, em conformidade aos conceitos embarcados descritos por \citeonline{white2024}. Um microcontrolador mestre centraliza as decisões matriciais analíticas e gerencia o envio assíncrono de comandos de geometria vetorial. Simultaneamente, um dispositivo escravo herda passivamente os segmentos de movimento oriundos do supervisor, atuando nas articulações periféricas do pulso. Deste modo, o paralelismo computacional preserva o tempo real crítico para o controle contínuo dos seis eixos do braço.

\subsection{O Robô Industrial de Referência (UR-10)}

No intuito de estipular balizas científicas validadas em âmbito produtivo escalável, delineou-se um comparativo laboratorial primário. O equipamento referencial escolhido foi o manipulador Universal Robots UR-10, consagrado amplamente na indústria por sua reprodutibilidade e conformidade normativa. Esse sistema dispõe de rastreabilidade completa em malha fechada restrita a desvios tolerantes abaixo de um milímetro real, conforme atestado por catálogos certificados \cite{universalrobots2023}. Consequentemente, a adoção do robô escandinavo permite aferir quantitativamente as limitações físicas do protótipo desenvolvido.

\section{Procedimentos de Implementação Tecnológica}

As resoluções de codificação eletrônicas objetivam forjar resiliência operacional contínua no braço robótico de baixo custo. Ao lidar com servomecanismos simplificados, isentos de respostas ativas embarcadas como \textit{encoders} absolutos integrados, a confiabilidade recai predominantemente na estruturação algorítmica. O planejamento lógico prevê contornar os gargalos de latência por meio de um sistema hierárquico imune a bloqueios sistêmicos e perda instantânea de pacotes. Dessa maneira, a seção aborda a consolidação das interações virtuais necessárias para a sincronia plena do dispositivo.

Por consequência, detalham-se as técnicas adotadas na programação compilada, capazes de emular a robustez de sistemas industriais na medida do possível. A arquitetura restaura a segurança de operação por meio de rotinas de vigilância acopladas aos ciclos repetitivos de \textit{software}. Adicionalmente, as instâncias de \textit{interface} proporcionam painéis manipuladores integrados nativamente na retaguarda através da rede sem fio local. Portanto, a implementação une a pragmática do controle digital de sinais com os requisitos de interface da automação remota.

\subsection{Mitigação Mecânica e Restauro de Posição}

Em dissonância com os padrões industriais estabelecidos, a matriz robótica inicial foi concebida desprovida de sensores perimetrais ou fim de curso. Diante dessa realidade, estipulou-se um procedimento preparatório no qual o operador assume manualmente a adequação das juntas perante marcadores físicos adesivos. Trata-se de uma calibração estática inicial que referencia permanentemente o plano relacional de cinemática para a sessão operacional ativa. Sendo assim, este eixo de coordenada zero serve de alicerce contínuo para o processamento trigonométrico subsequente.

Todavia, para assegurar a coerência sistêmica contra interrupções desastrosas ou paradas abruptas, as posições calculadas são salvas periodicamente. Engajam-se rotinas explícitas que flanqueiam cada movimento longo com persistência de dados em memórias não voláteis integradas ao controlador principal. Dessa feita, em todo evento de religamento sistêmico, os estados vetoriais elétricos decodificados renascem de forma idêntica via leitura isolada da memória de armazenamento local. Consequentemente, mantêm-se ativos e robustos os limites virtuais de segurança das engrenagens impressas.

\subsection{Protocolo de Enlace Mestre-Escravo}

Sincronizar espacialmente os dois microcontroladores obrigou o desenvolvimento de pontes de controle serial estritas e determinísticas. O formato imposto consistiu em um canal UART isolado com arquitetura de pacotes limitados por cabeçalhos e finalizadores constantes. Incorporaram-se atestados de integridade baseados em verificações cíclicas de redundância aos vetores de comando, descartando ordens possivelmente malformadas por interferência indutiva. Por conseguinte, as lógicas do supervisor não saturam passivamente a capacidade de processamento do módulo escravo periférico.

Ademais, essa malha abrange mecanismos de devoluções interativas de conformidade ancoradas no eco afirmativo entre os processadores. O escravo pulsa ativamente sinais de garantia atestando a finalização completa das metas de movimentação comandadas instantes antes pelo mestre. Falhas contínuas de acoplamento acionam rotinas emergenciais que desativam sumariamente o fornecimento de potência para a matriz de motores em ambas as extremidades. Desse modo, previnem-se danos estruturais no protótipo derivados de movimentos caóticos assíncronos não supervisionados.

\subsection{Interface e Interconectividade}

Pautado na exploração conceitual da Internet das Coisas (\textit{IoT}) aplicada à fronteira experimental, implementou-se um serviço de rádio local exclusivo. A manobra restringe as instabilidades inerentes às conexões externas de instituições acadêmicas, ao priorizar o modo ponto de acesso direto hospedado no próprio mestre. Esse escopo isolado irradia um painel de supervisão dinâmico capaz de receber interações cruciais do operador de maneira responsiva. Portanto, o ambiente assegura uma via de teste ininterrupta sem dependência de plataformas de nuvem instáveis.

\section{Protocolos Experimentais e Configuração de Testes}

A validação paramétrica em escopo prático constitui o referencial de aprovação das rotinas desenvolvidas perante os objetivos propostos. Respaldando-se ativamente por orientações normativas consolidadas, como a ISO 9283, as medições isolam fatores de desvio de repetibilidade em cenários delimitados de atuação forçada \citeonline{iso9283}. A sequência de testes transcende os algoritmos de bordo, avaliando o conjunto na execução de tarefas tridimensionais similares aos corriqueiros manuseios fabris. Desse modo, afere-se a confiabilidade real de longo prazo do equipamento sob estresse simulado contínuo.

Dessarte, a consecução dessas diretrizes materializa-se através da fragmentação empírica dos módulos eletromecânicos acoplados. Contornam-se as barreiras orçamentárias de aferição visual computadorizada avançada pelo emprego de recursos óticos passivos calibrados em bancada simples de trabalho acadêmico. Adicionalmente, as interações lógicas são estressadas à exaustão usando clientes artificiais que disparam comandos seriais massivos para testar a resiliência comunicacional da interface. Sendo assim, o arcabouço comprobatório cruza fatores físicos, lógicos e temporais com eficácia.

\subsection{Ensaios Mecânicos e Lógicos Estruturados}

A fase operacional da construção exige rotinas cadenciadas, focando na integridade das engrenagens impressas sob estresse dinâmico progressivo. Inicialmente, impõem-se torções independentes limitadas em cada eixo tracional para observar e registrar defasagens inerciais inerentes aos engates plásticos submetidos a peso fixo. As articulações operam em simulações isoladas até atingir o aquecimento estabilizado dos blocos estatores perante atrito controlável e continuado. Com base nisso, os perfis mecânicos das fraquezas estruturais norteiam as compensações introduzidas futuramente no código-fonte matricial principal.

Em contrapartida, as baterias de ensaios lógicos analisam o estrangulamento de pacotes através do esgotamento intencional dos buffers sistêmicos. Satura-se o servidor embarcado com requisições oriundas da rede remota local enquanto operações críticas de posicionamento motorizado estão ocorrendo. A eficácia dos sistemas de segurança e vigilância algorítmica prova-se ao desativar ou mitigar as anomalias geradas por comandos corrompidos injetados na malha. Assim, isola-se o risco elétrico derivado puramente das flutuações de programação do dispositivo mestre inteligente.

\subsection{Qualificação Funcional da Resposta Prática}

O desafio analítico tangibiliza as comparações de precisão adotando mapeamentos visuais contínuos e medições mecânicas em um plano quadriculado. Efetua-se marcações no perímetro de operação abrangendo raios de alcance práticos para validar tempos de latência e desvios de repetição milimétrica nos alvos coordenados programados. Tais dinâmicas compensam a ausência de sistemas complexos a laser, mantendo a documentação e cronometria fiel por meio de medições passivas assistidas sistematicamente por registro óptico e tátil minucioso. Dessa forma, as discrepâncias acumuladas pelo escorregamento microscópico de eixos transparecem quantitativamente.

Ademais, as rotinas estabelecem ciclos ininterruptos prolongados para forçar e medir as metas ditadas pela normativa de repetibilidade espacial. As varreduras mapeiam a perda de posição em virtude dos micropassos motores esquipados acidentalmente frente à aceleração estrita demandada em cada vértice da jornada virtual. O equipamento é avaliado quanto à persistência de parada e ao esgotamento por fadiga das interações seriais nos comandos de malha aberta restrita. Portanto, certifica-se o envelope útil viável do protótipo dentro das aplicações didáticas ou simplificadas propostas.

\subsection{Emulação de Aplicações Industriais Reais}

Como método prático definitivo de demonstração das capacidades sistêmicas, consolida-se a implementação utilitarista na forma de movimentações sintéticas corriqueiras do cotidiano fabril. Propõe-se um protocolo dinâmico prolongado de exaustividade nomeado \textit{Pick and Place}, cujo intuito é transportar repetitivamente massas volumétricas padronizadas. Os ensaios analisam desvios incrementais da garra extratora em longas simulações térmicas sujeitas a perdas interativas nas posições calculadas estaticamente no espaço de atuação. Desse modo, a simulação emula fidedignamente o stress esperado num ambiente de uso automatizado contínuo.

Além disso, etapas graduais introduzem fluxos de comunicação assíncronos e remotos para compor a tomada de decisão do equipamento central durante o transporte. As simulações logísticas integram envios via protocolo HTTP que engatilham a série motora autônoma, validando integralmente o conceito de automação baseada na Internet das Coisas (\textit{IoT}). Consequentemente, as interações ativas validam a resposta temporal do braço articulado robótico no escopo pretendido inicialmente pela pesquisa de engenharia, demonstrando êxito empírico e técnico acadêmico funcional.

